\section{Immutable Parent Pointers}

\subsection{Motivation for Immutability}

In any recursive spawning architecture, the \textbf{parent-child relationship} constitutes the foundational metadata that governs agent lineage, accountability, and causal tracing. Once a spawn event assigns a parent to a newly created agent, that parent pointer must remain invariant throughout the agent's lifetime. Mutable parent pointers would introduce ambiguity in causal chains, compromise auditability, and undermine the very premise of recursive accountability that the architecture seeks to enforce.

Consider the failure modes that immutability prevents. If a parent pointer could be retrospectively modified, lineage records would become unreliable: an agent's provenance could change retroactively, breaking invariants that downstream systems depend upon.溯源 (traceability) would be compromised. Authorization decisions made at spawn time—predicated on the parent pointer's value—could be invalidated post-hoc. The Fact Layer's record of what \emph{was} would diverge from the runtime state of what \emph{is}.

The immutability constraint therefore serves as a \textbf{structural guarantee} rather than a mere implementation preference. It ensures that the spawn authorization recorded by the Purpose Layer remains meaningfully coupled to the factual record maintained by the Fact Layer. The design enforces that once a parent is assigned, it is assigned permanently, with no mechanism provided for subsequent modification.\footnote{This design choice aligns with append-only data architectures in distributed systems, where immutability enables reasoning about history without requiring locking or consensus on past events.}

\subsection{Formal Definition: ParentPointer Relation}

We formalize the parent-child relationship through the binary relation \textbf{ParentPointer}.

\begin{definition}[ParentPointer]
Let $p, c \in \mathcal{A}$ denote two agents in the agent population $\mathcal{A}$. The relation $\ParentPointer(p, c)$ holds \textbf{iff} $p$ is the sole, immutable parent of $c$ at the moment of $c$'s creation. Formally:
\begin{equation}
\ParentPointer(p, c) \iff \text{spawn}(c) \text{ was authorized by } p \land \text{no subsequent reassignment is permitted.}
\end{equation}
\end{definition}

The ParentPointer relation is established \textbf{atomically} at the moment of spawn authorization. The Purpose Layer's authorization event causes the Fact Layer to record the assignment. There is no temporal window in which the relation is undefined, and there is no subsequent state transition that can alter it.

Key properties of $\ParentPointer$:

\begin{itemize}
    \item \textbf{Uniqueness:} For any agent $c \neq \text{null}$, there exists exactly one $p$ such that $\ParentPointer(p, c)$ holds.
    \item \textbf{Immutability:} Once $\ParentPointer(p, c)$ is recorded, no operation in any layer may modify or revoke it.
    \item \textbf{Atomic assignment:} The relation is established in a single, atomic Fact Layer write during the spawn event.
\end{itemize}

We represent the parent of an agent $a$ formally as a function:

\begin{equation}
parent(a) = \begin{cases}
p & \text{if } \ParentPointer(p, a) \\
\text{null} & \text{if } a \text{ is a root human}
\end{cases}
\end{equation}

The $\parent$ function is total over $\mathcal{A}$; every agent has a defined parent value, either an agent reference or $\text{null}$ for root humans.

\subsection{Recursive Chain Definition}

Given the $\parent$ function, we define the \textbf{ancestral chain} of an agent recursively. The chain encodes the complete lineage from the agent back to the root human, encompassing all intermediate spawning relationships.

\begin{definition}[Ancestral Chain]
For any agent $a \in \mathcal{A}$, the ancestral chain $\Chain(a)$ is defined recursively as:
\begin{equation}
\Chain(a) = \ParentPointer(\parent(a), a) \cup \Chain(\parent(a))
\end{equation}
with base case:
\begin{equation}
\Chain(r) = \emptyset \quad \text{where } r \text{ is a root human and } \parent(r) = \text{null}.
\end{equation}
\end{definition}

Expanding the recursive definition, we can express $\Chain(a)$ explicitly as the set:
\begin{equation}
\Chain(a) = \{ \ParentPointer(\parent(a), a), \ParentPointer(\parent(\parent(a)), \parent(a)), \ParentPointer(\parent(\parent(\parent(a))), \parent(\parent(a))), \ldots \}
\end{equation}

This chain terminates when reaching the root human $r$, at which point the union with $\Chain(r) = \emptyset$ ends the recursion.

An equivalent iterative definition clarifies implementation:
\begin{equation}
\Chain(a) = \{ \ParentPointer(\parent^i(a), \parent^{i-1}(a)) \mid i \in \{1, 2, \ldots, n\} \}
\end{equation}
where $\parent^n(a)$ is the $n$-fold application of $\parent$, and $\parent^n(a) = \text{null}$ (i.e., $a$'s $n$-th ancestor is the root human).

\begin{figure}[tb]
\centering
\fbox{
\begin{minipage}{0.92\linewidth}
\begin{center}
\textbf{Chain Construction Illustration}
\end{center}
\vspace{6pt}
Suppose we have agents $A$, $B$, $C$, $D$ with spawn relationships:
\begin{align*}
\ParentPointer(H, A) &\quad \text{(A spawned by human H, root)} \\
\ParentPointer(A, B) &\quad \text{(B spawned by A)} \\
\ParentPointer(B, C) &\quad \text{(C spawned by B)} \\
\ParentPointer(C, D) &\quad \text{(D spawned by C)}
\end{align*}
Then:
\begin{align*}
\Chain(D) &= \{ \ParentPointer(C, D) \} \cup \Chain(C) \\
         &= \{ \ParentPointer(C, D), \ParentPointer(B, C) \} \cup \Chain(B) \\
         &= \{ \ParentPointer(C, D), \ParentPointer(B, C), \ParentPointer(A, B) \} \cup \Chain(A) \\
         &= \{ \ParentPointer(C, D), \ParentPointer(B, C), \ParentPointer(A, B) \} \cup \emptyset \\
         &= \{ \ParentPointer(C, D), \ParentPointer(B, C), \ParentPointer(A, B) \}
\end{align*}
\end{minipage}
}
\caption{Recursive chain construction traces all ancestors back to root human.}
\label{fig:chain-construction}
\end{figure}

The chain encodes the complete spawn genealogy. It is inherently acyclic by construction: the recursive definition ensures no agent can appear as its own ancestor, as spawn relationships are established only from existing agents to new agents.

\subsection{The Root Human as Terminal Base Case}

The base case of the chain recursion is the \textbf{root human}. A root human is a human user who exists prior to and independent of the AgentOS architecture. They are not themselves agents spawned by the system; they are the originating principals who initiate the first layer of agent creation.

Formally:

\begin{align*}
\text{RootHumans} &= \{ h \in \mathcal{H} \mid \nexists a \in \mathcal{A} : \ParentPointer(a, h) \} \\
\text{RootHuman}(h) &\iff h \in \text{RootHumans}
\end{align*}

where $\mathcal{H}$ denotes the set of human users within the system.

For a root human $r$:
\begin{equation}
\parent(r) = \text{null}
\end{equation}
and therefore:
\begin{equation}
\Chain(r) = \emptyset.
\end{equation}

The null parent value for root humans is a critical design choice. It establishes a well-defined termination point for chain traversal. Without a null marker, recursion could not terminate cleanly—the chain would require a sentinel value to signal termination, and the recursive definition would remain incomplete.

Note that $\text{null}$ is not itself an agent; it is a distinguished value indicating the absence of a parent. The set of agents $\mathcal{A}$ contains only non-null entities. Root humans, while they are humans, are not agents in the spawn graph—they are external principals that stand as origins.

Implications of this design:

\begin{itemize}
    \item Every agent has a finite chain terminating at a root human.
    \item Chain traversal algorithms are guaranteed to terminate.
    \item溯源 queries can always reach a human principal by following the chain.
    \item The root human provides the ultimate accountability anchor.
\end{itemize}

\subsection{Immutability Enforcement: Fact Layer Architecture}

The Fact Layer provides the \textbf{enforcement mechanism} for parent pointer immutability. The Fact Layer is the authoritative record of all factual state in the AgentOS architecture, including spawn events, parent assignments, and chain relationships.

Immutability is not merely a convention—it is architecturally enforced:

\begin{enumerate}
    \item \textbf{Write-once assignment.} When the Purpose Layer authorizes a spawn event, the Fact Layer records the parent pointer as an immutable attribute of the child agent. The Fact Layer provides no API operation for modifying an existing parent pointer. There is no $\text{setParent}(a, p')$ function—only the initial assignment at spawn time.
    
    \item \textbf{Integrity constraints.} The Fact Layer schema enforces the uniqueness constraint: a non-root agent must have exactly one parent. Any write operation that would create a second parent assignment for an existing agent is rejected at the schema level. This is implemented as a NOT NULL constraint on the parent reference and a uniqueness constraint on the child side of the relation.

    \item \textbf{Append-only event log.} The underlying event log of the Fact Layer is append-only. Spawn events are recorded as new entries; there is no entry type for parent reassignment. The log is structurally incapable of representing a modification to an existing parent pointer.

    \item \textbf{Audit trail.} Every Fact Layer write is timestamped and immutable. The record of who authorized the spawn (Purpose Layer), what parent was assigned, and at what time constitutes the authoritative lineage record. There is no mechanism within the Fact Layer to alter this record after the fact.
\end{enumerate}

\begin{figure}[tb]
\centering
\fbox{
\begin{minipage}{0.92\linewidth}
\begin{center}
\textbf{Fact Layer Enforcement of ParentPointer Immutability}
\end{center}
\vspace{6pt}
\begin{tabular}{|l|l|}
\hline
\textbf{Operation} & \textbf{Fact Layer Response} \\
\hline
Spawn with parent $p$ assigned & Record $\ParentPointer(p, c)$; timestamp; commit \\
\hline
Attempt to reassign parent($c$) = $p'$ & Reject; no such operation exists \\
\hline
Attempt to delete $\ParentPointer(p, c)$ & Reject; relation is delete-protected \\
\hline
Query parent($c$) after spawn & Return $p$ (immutable) \\
\hline
Query Chain($c$) & Return complete chain to root human \\
\hline
\end{tabular}
\end{minipage}
}
\caption{Fact Layer enforces immutability through schema and API design.}
\label{fig:fact-layer-immutability}
\end{figure}

The Fact Layer thus provides a \textbf{tamper-evident} record. Any attempt to circumvent the immutability constraint would require modifying the Fact Layer's append-only log, which is protected by the system's operational security controls. The architecture assumes that the Fact Layer itself is immutable unless the system itself is compromised, at which point the guarantees of the architecture are by definition no longer operative.

\subsection{Spawn Authorization: Purpose Layer Constraints}

The Purpose Layer governs \textbf{why} agents may be created. The Purpose Layer authorization must precede the Fact Layer recording of the parent pointer. The authorization identifies the parent as a prerequisite for the Fact Layer to record the $\ParentPointer$ relation.

The spawn authorization process:

\begin{enumerate}
    \item A principal (human or agent) requests to spawn a new agent $c$.
    \item The Purpose Layer evaluates whether the principal is authorized to spawn.
    \item If authorized, the Purpose Layer records the spawn intent, including the designated parent $p$.
    \item The Fact Layer receives the authorization and records $\ParentPointer(p, c)$ as an immutable attribute of $c$.
    \item The parent-child relationship is now fixed and permanent.
\end{enumerate}

The Purpose Layer thus acts as the \textbf{gateway} to parent pointer assignment. Without valid authorization, no parent pointer is recorded. With valid authorization, the parent pointer is recorded and locked. There is no mechanism to reassign the parent after authorization—the authorization event itself triggers the immutable recording.

This design ensures that the parent pointer always reflects an authorized spawning decision made at spawn time, not a retrospective post-hoc assignment. The Purpose Layer's decision is the proximate cause of the parent pointer value, and the immutability of the Fact Layer record ensures that this causal link is preserved indefinitely.

The Purpose Layer may impose constraints on valid parent choices. For instance, it may require that the parent be an active agent with sufficient privileges, or that the spawn request passes certain policy checks. These constraints are evaluated at authorization time and encoded in the Purpose Layer record. The Fact Layer records the outcome—not the reasoning—ensuring separation of concerns between authorization logic and factual record-keeping.

\subsection{Chain Traversal Algorithm}

Chain traversal is required for溯源 queries, authorization checks, and audit operations. The algorithm must correctly handle the recursive structure, terminate at the root human, and respect immutability guarantees (i.e., it reads only, never writes).

\begin{algorithm}[tb]
\caption{Chain Traversal: Compute Ancestral Chain}
\label{algo:chain-traversal}
\begin{algorithmic}
\Procedure{GetChain}{$a$}
    \State $\text{chain} \gets \emptyset$
    \State $\text{current} \gets a$
    \While{$\text{current} \neq \text{null}$}
        \State $p \gets \parent(\text{current})$
        \If{$p = \text{null}$}
            \State \Comment{Root human reached; chain terminates}
            \State $\textbf{break}$
        \EndIf
        \State $\text{chain} \gets \text{chain} \cup \{ \ParentPointer(p, \text{current}) \}$
        \State $\text{current} \gets p$
    \EndWhile
    \State \Return $\text{chain}$
\EndProcedure
\end{algorithmic}
\end{algorithm}

The iterative algorithm correctly implements the recursive definition of $\Chain(a)$. The loop invariant is that at each iteration, $\text{current}$ holds the next ancestor to be added to the chain, and $\text{chain}$ holds all ancestors processed so far. The loop terminates when $\text{current}$ becomes $\text{null}$, indicating the root human has been reached and the chain is complete.

\textbf{Correctness argument:} We prove that $\text{GetChain}(a)$ returns $\Chain(a)$ as defined in Definition 2. We proceed by induction on the number of iterations $k$ of the while loop.

\begin{itemize}
    \item \textbf{Base case ($k = 0$):} Before the first iteration, $\text{current} = a$ and $\text{chain} = \emptyset$. The chain is empty, and no ancestors have been processed. This matches the initial state of the recursive definition before any union operations.
    
    \item \textbf{Inductive step:} Assume after $k$ iterations that $\text{chain}$ contains $\{ \ParentPointer(\parent^i(a), \parent^{i-1}(a)) \mid i \in \{1, \ldots, k\} \}$ and $\text{current} = \parent^k(a)$. On iteration $k+1$, we retrieve $p = \parent(\text{current}) = \parent^{k+1}(a)$. If $p \neq \text{null}$, we add $\ParentPointer(p, \text{current}) = \ParentPointer(\parent^{k+1}(a), \parent^k(a))$ to $\text{chain}$ and set $\text{current} = p$. Thus $\text{chain}$ now contains all ancestors up to $\parent^k(a)$, and $\text{current}$ holds $\parent^{k+1}(a)$.
    
    \item \textbf{Termination:} The loop terminates when $p = \text{null}$, meaning $\parent^{n+1}(a) = \text{null}$ for some $n$. This occurs precisely when $\parent^n(a)$ is a root human (its parent is null). At this point, $\text{chain}$ contains all edges from $\ParentPointer(\parent(a), a)$ through $\ParentPointer(\parent^n(a), \parent^{n-1}(a))$, which by definition is $\Chain(a)$. The algorithm then returns $\text{chain}$.
\end{itemize}

\textbf{Complexity analysis:} Each iteration removes one ancestor from the chain (advances toward the root). Since every agent has a finite chain length bounded by the depth of the spawn tree, the algorithm terminates after at most $O(d)$ iterations, where $d$ is the maximum chain depth. The time complexity is $O(d)$ and the space complexity is $O(d)$ for storing the resulting chain.

An important variant is the \textbf{ancestor query}: checking whether a given agent $x$ appears in the chain of another agent $a$. This can be implemented efficiently by traversing the chain and comparing each ancestor to $x$, terminating early upon match or when the root is reached:

\begin{algorithm}[H]
\caption{Ancestor Check: Is $x$ an ancestor of $a$?}
\label{algo:ancestor-check}
\begin{algorithmic}
\Procedure{IsAncestor}{$a, x$}
    \State $\text{current} \gets a$
    \While{$\text{current} \neq \text{null}$}
        \State $p \gets \parent(\text{current})$
        \If{$p = x$}
            \State \Return $\text{true}$
        \EndIf
        \State $\text{current} \gets p$
    \EndWhile
    \State \Return $\text{false}$
\EndProcedure
\end{algorithmic}

This algorithm also runs in $O(d)$ time and is optimal for single-query scenarios. For bulk queries, indexing structures (e.g., ancestor tables) could be maintained as derived artifacts of the Fact Layer, but the fundamental traversal remains correct.

\begin{figure}[tb]
\centering
\fbox{
\begin{minipage}{0.92\linewidth}
\begin{center}
\textbf{Example: Chain Traversal Queries}
\end{center}
\vspace{6pt}
Given the spawn graph from Figure~\ref{fig:chain-construction}:
\begin{align*}
\ParentPointer(H, A), \ParentPointer(A, B), \ParentPointer(B, C), \ParentPointer(C, D)
\end{align*}
\begin{tabular}{|l|l|l|}
\hline
\textbf{Query} & \textbf{Result} & \textbf{Explanation} \\
\hline
$\Chain(D)$ & $\{ \ParentPointer(C,D), \ParentPointer(B,C), \ParentPointer(A,B) \}$ & D's full lineage to root H \\
\hline
$\Chain(B)$ & $\{ \ParentPointer(A,B) \}$ & B's parent is A; A's parent is null (root) \\
\hline
$\Chain(H)$ & $\emptyset$ & H is a root human; null parent terminates immediately \\
\hline
$\IsAncestor(D, B)$ & true & B appears as $C$'s parent in D's chain \\
\hline
$\IsAncestor(D, H)$ & true & H appears as A's parent in D's chain (root) \\
\hline
$\IsAncestor(B, D)$ & false & D is a descendant of B, not an ancestor \\
\hline
$\IsAncestor(C, A)$ & true & A is a grandparent of C (appears in chain) \\
\hline
\end{tabular}
\end{minipage}
}
\caption{Chain traversal queries on a four-agent spawn tree.}
\label{fig:chain-queries}
\end{figure}

\subsection{Properties and Invariants}

The combination of immutable parent pointers, recursive chain definition, and Fact Layer enforcement yields several important invariants that hold throughout the system.

\begin{prop}[Finite Chain Property]
For any agent $a \in \mathcal{A}$, $\Chain(a)$ is finite and terminates at a root human.
\end{prop}
\begin{Proof}
The spawn graph is acyclic: agents are only created from existing agents or from human principals. There is no operation that creates a cycle in the parent relationship. Since every spawn creates a new agent with a parent drawn from the existing population, the chain cannot be infinite. The only terminal node is a root human (whose parent is null), and all chains must eventually reach null. $\square$
\end{Proof}

\begin{prop}[Uniqueness of Lineage]
For any two agents $a, b \in \mathcal{A}$ with $a \neq b$, either $\Chain(a) \cap \Chain(b) = \emptyset$ or one chain is a strict subset of the other (ancestor-descendant relationship).
\end{prop}
\begin{Proof}
Two distinct agents cannot share the same immediate parent (uniqueness constraint). By induction on chain depth, they cannot share any ancestor other than a common ancestor higher in the tree. If they share a common ancestor, one must be descended from the other, yielding a subset relationship. If they have different root humans, their chains are disjoint. $\square$
\end{Proof}

\begin{prop}[Authorization-Record Coupling]
If $\ParentPointer(p, c)$ is recorded in the Fact Layer, then the Purpose Layer authorized the spawn of $c$ by $p$ at that time.
\end{prop}
\begin{Proof}
The Fact Layer records parent pointers only upon receiving a valid authorization from the Purpose Layer. There is no Fact Layer API to record a parent pointer without a corresponding authorization event. Therefore, every recorded $\ParentPointer$ relation has a corresponding authorized spawn event. $\square$
\end{Proof}

\begin{prop}[Immutability Soundness]
No operation within the system's normal operational modes can alter a recorded $\ParentPointer$ relation.
\end{prop}
\begin{Proof}
The Fact Layer provides no mutation API for parent pointers (no update, no reassign, no delete). The schema enforces uniqueness and non-null constraints. The event log is append-only. Any attempt to modify the parent pointer would require a mechanism that does not exist in the Fact Layer's operational interface. $\square$
\end{Proof}

These invariants collectively ensure that the spawn graph is a \textbf{valid, immutable, auditable tree structure} rooted at human principals. The immutability of parent pointers is not a policy that can be relaxed—it is enforced by the architecture, making the tree structure a fundamental property of the system rather than a convention that could be violated.

\subsection{Summary}

The \textbf{Immutable Parent Pointer} design is central to the recursive spawning architecture's accountability model. The relation $\ParentPointer(p, c)$ is established atomically at spawn time, is unique for each child, and is permanently recorded in the Fact Layer with no mechanism for modification. The recursive definition $\Chain(a) = \ParentPointer(\parent(a), a) \cup \Chain(\parent(a))$ builds the complete ancestral lineage from any agent back to the root human, with the base case being the null parent of a root human.

Immutability is architecturally enforced by the Fact Layer through write-once assignment, schema constraints, append-only logging, and audit trails. The Purpose Layer gates spawn authorization, ensuring that parent assignments arise from authorized decisions made at spawn time. Chain traversal algorithms—iterative or recursive—correctly compute ancestral chains and terminate at the root human due to the finite chain property.

These mechanisms together establish the parent pointer as an immutable, trustworthy foundation for溯源, authorization, and audit operations throughout the AgentOS architecture. The design ensures that every agent's lineage is permanently recorded, permanently immutable, and permanently traceable to a human principal.